\section{Chow cohomology and the first Chern class}
\label{section-chow-cohomology}

\noindent
We will be most interested in $A^p(X) = A^p(X \to X)$, which will always mean
the bivariant cohomology classes for $\text{id}_X$. Namely, that is where
Chern classes will live.

\begin{definition}
\label{definition-chow-cohomology}
Let $(S, \delta)$ be as in Situation \ref{situation-setup}.
Let $X$ be locally of finite type over $S$. The {\it Chow cohomology}
of $X$ is the graded $\mathbf{Z}$-algebra $A^*(X)$ whose degree
$p$ component is $A^p(X \to X)$.
\end{definition}

\noindent
Warning: It is not clear that the $\mathbf{Z}$-algebra structure
on $A^*(X)$ is commutative, but we will see that Chern classes live
in its center.

\begin{remark}
\label{remark-pullback-cohomology}
Let $(S, \delta)$ be as in Situation \ref{situation-setup}.
Let $f : Y' \to Y$ be a morphism of schemes locally of finite type over $S$.
As a special case of Remark \ref{remark-restriction-bivariant}
there is a canonical $\mathbf{Z}$-algebra map $res : A^*(Y) \to A^*(Y')$.
This map is often denoted $f^*$ in the literature.
\end{remark}

\begin{lemma}
\label{lemma-cap-c1-bivariant}
Let $(S, \delta)$ be as in Situation \ref{situation-setup}.
Let $X$ be locally of finite type over $S$.
Let $\mathcal{L}$ be an invertible $\mathcal{O}_X$-module.
Then the rule that to $f : X' \to X$ assigns
$c_1(f^*\mathcal{L}) \cap - : \CH_k(X') \to \CH_{k - 1}(X')$
is a bivariant class of degree $1$.
\end{lemma}

\begin{proof}
This follows from Lemmas \ref{lemma-factors},
\ref{lemma-pushforward-cap-c1},
\ref{lemma-flat-pullback-cap-c1}, and
\ref{lemma-gysin-commutes-cap-c1}.
\end{proof}

\noindent
The lemma above finally allows us to make the following definition.

\begin{definition}
\label{definition-first-chern-class}
Let $(S, \delta)$ be as in Situation \ref{situation-setup}. Let $X$
be locally of finite type over $S$. Let $\mathcal{L}$ be an invertible
$\mathcal{O}_X$-module. The {\it first Chern class}
$c_1(\mathcal{L}) \in A^1(X)$ of $\mathcal{L}$
is the bivariant class of Lemma \ref{lemma-cap-c1-bivariant}.
\end{definition}

\noindent
For finite locally free modules we construct the Chern classes in
Section \ref{section-intersecting-chern-classes}.
Let us prove that $c_1(\mathcal{L})$ is in the center of $A^*(X)$.

\begin{lemma}
\label{lemma-c1-center}
Let $(S, \delta)$ be as in Situation \ref{situation-setup}.
Let $X$ be locally of finite type over $S$.
Let $\mathcal{L}$ be an invertible $\mathcal{O}_X$-module.
Then
\begin{enumerate}
\item $c_1(\mathcal{L}) \in A^1(X)$ is in the center of $A^*(X)$ and
\item if $f : X' \to X$ is locally of finite type and $c \in A^*(X' \to X)$,
then $c \circ c_1(\mathcal{L}) = c_1(f^*\mathcal{L}) \circ c$.
\end{enumerate}
\end{lemma}

\begin{proof}
Of course (2) implies (1).
Let $p : L \to X$ be as in Lemma \ref{lemma-linebundle} and let $o : X \to L$
be the zero section. Denote $p' : L' \to X'$ and $o' : X' \to L'$
their base changes. By Lemma \ref{lemma-linebundle-formulae} we have
$$
p^*(c_1(\mathcal{L}) \cap \alpha) = - o_* \alpha
\quad\text{and}\quad
(p')^*(c_1(f^*\mathcal{L}) \cap \alpha') = - o'_* \alpha'
$$
Since $c$ is a bivariant class we have
\begin{align*}
(p')^*(c \cap c_1(\mathcal{L}) \cap \alpha)
& =
c \cap p^*(c_1(\mathcal{L}) \cap \alpha) \\
& =
- c \cap o_* \alpha \\
& =
- o'_*(c \cap \alpha) \\
& =
(p')^*(c_1(f^*\mathcal{L}) \cap c \cap \alpha)
\end{align*}
Since $(p')^*$ is injective by one of the lemmas cited above we obtain
$c \cap c_1(\mathcal{L}) \cap \alpha =
c_1(f^*\mathcal{L}) \cap c \cap \alpha$.
The same is true after any base change by $Y \to X$ locally of finite type
and hence we have the equality of bivariant classes stated in (2).
\end{proof}

\begin{lemma}
\label{lemma-vanish-above-dimension}
Let $(S, \delta)$ be as in Situation \ref{situation-setup}. Let $X$ be a
finite type scheme over $S$ which has an ample invertible sheaf.
Assume $d = \dim(X) < \infty$ (here we really mean dimension and
not $\delta$-dimension).
Then for any invertible sheaves $\mathcal{L}_1, \ldots, \mathcal{L}_{d + 1}$
on $X$ we have
$c_1(\mathcal{L}_1) \circ \ldots \circ c_1(\mathcal{L}_{d + 1}) = 0$
in $A^{d + 1}(X)$.
\end{lemma}

\begin{proof}
We prove this by induction on $d$. The base case $d = 0$ is true because
in this case $X$ is a finite set of closed points and hence every invertible
module is trivial. Assume $d > 0$. By Divisors, Lemma
\ref{divisors-lemma-quasi-projective-Noetherian-pic-effective-Cartier}
we can write $\mathcal{L}_{d + 1} \cong \mathcal{O}_X(D) \otimes
\mathcal{O}_X(D')^{\otimes -1}$ for some effective Cartier divisors
$D, D' \subset X$. Then $c_1(\mathcal{L}_{d + 1})$ is the difference
of $c_1(\mathcal{O}_X(D))$ and $c_1(\mathcal{O}_X(D'))$ and hence
we may assume $\mathcal{L}_{d + 1} = \mathcal{O}_X(D)$ for some
effective Cartier divisor.

\medskip\noindent
Denote $i : D \to X$ the inclusion morphism and denote
$i^* \in A^1(D \to X)$ the bivariant class given by the
Gysin hommomorphism as in Lemma \ref{lemma-gysin-bivariant}.
We have $i_* \circ i^* = c_1(\mathcal{L}_{d + 1})$
in $A^1(X)$ by Lemma \ref{lemma-support-cap-effective-Cartier}
(and Lemma \ref{lemma-push-proper-bivariant}
to make sense of the left hand side).
Since $c_1(\mathcal{L}_i)$ commutes with
both $i_*$ and $i^*$ (by definition of bivariant classes)
we conclude that
$$
c_1(\mathcal{L}_1) \circ \ldots \circ c_1(\mathcal{L}_{d + 1}) =
i_* \circ c_1(\mathcal{L}_1) \circ \ldots \circ c_1(\mathcal{L}_d) \circ i^* =
i_* \circ c_1(\mathcal{L}_1|_D) \circ \ldots \circ c_1(\mathcal{L}_d|_D)
\circ i^*
$$
Thus we conclude by induction on $d$. Namely, we have $\dim(D) < d$
as none of the generic points of $X$ are in $D$.
\end{proof}

\begin{remark}
\label{remark-ring-loc-classes}
Let $(S, \delta)$ be as in Situation \ref{situation-setup}.
Let $Z \to X$ be a closed immersion of schemes locally of
finite type over $S$ and let $p \geq 0$. In this setting we define
$$
A^{(p)}(Z \to X) =
\prod\nolimits_{i \leq p - 1} A^i(X) \times
\prod\nolimits_{i \geq p} A^i(Z \to X).
$$
Then $A^{(p)}(Z \to X)$ canonically comes equipped with the structure
of a graded algebra. In fact, more generally there is a multiplication
$$
A^{(p)}(Z \to X) \times A^{(q)}(Z \to X)
\longrightarrow A^{(\max(p, q))}(Z \to X)
$$
In order to define these we define maps
\begin{align*}
A^i(Z \to X) \times A^j(X) & \to A^{i + j}(Z \to X) \\
A^i(X) \times A^j(Z \to X) & \to A^{i + j}(Z \to X) \\
A^i(Z \to X) \times A^j(Z \to X) & \to A^{i + j}(Z \to X)
\end{align*}
For the first we use composition of bivariant classes.
For the second we use restriction
$A^i(X) \to A^i(Z)$ (Remark \ref{remark-restriction-bivariant}) and
composition $A^i(Z) \times A^j(Z \to X) \to A^{i + j}(Z \to X)$.
For the third, we send $(c, c')$ to $res(c) \circ c'$
where $res : A^i(Z \to X) \to A^i(Z)$ is the restriction map (see
Remark \ref{remark-restriction-bivariant}). We omit the
verification that these multiplications are associative in a suitable sense.
\end{remark}

\begin{remark}
\label{remark-res-push}
Let $(S, \delta)$ be as in Situation \ref{situation-setup}.
Let $Z \to X$ be a closed immersion of schemes locally of
finite type over $S$. Denote $res : A^p(Z \to X) \to A^p(Z)$
the restriction map of Remark \ref{remark-restriction-bivariant}.
For $c \in A^p(Z \to X)$ we have
$res(c) \cap \alpha = c \cap i_*\alpha$ for $\alpha \in \CH_*(Z)$.
Namely $res(c) \cap \alpha = c \cap \alpha$
and compatibility of $c$ with proper pushforward
gives $(Z \to Z)_*(c \cap \alpha) = c \cap (Z \to X)_*\alpha$.
\end{remark}






